\documentclass[../../main]{subfiles}

\begin{document}
\section{Fundamentos}
\label{sec:matematica:logica-matematica:fundamentos}

\begin{defn}[Conceito Primitivo]
	\label{def:matematica:logica-matematica:fundamentos:conceito-primitivo}

	Um conceito primitivo é um conceito adotado sem definição formal dentro
	de uma teoria, servindo como fundamento para a construção de definições,
	axiomas e proposições.
\end{defn}

\begin{defn}[Proposição]
	\label{def:matematica:logica-matematica:fundamentos:proposicao}

	Uma proposição é uma sentença declarativa à qual pode ser atribuído
	exatamente um valor de verdade: verdadeiro ou falso.
\end{defn}

\begin{defn}[Axioma]
	\label{def:matematica:logica-matematica:fundamentos:axioma}

	Um axioma é uma proposição admitida sem demonstração dentro de um
	sistema axiomático.
\end{defn}

\begin{defn}[Sistema Axiomático]
	\label{def:matematica:logica-matematica:fundamentos:sistema-axiomatico}

	Um sistema axiomático é uma estrutura formal constituída por conceitos
	primitivos, axiomas e regras de inferência, a partir dos quais são
	definidas e demonstradas proposições.
\end{defn}

\begin{defn}[Demonstração]
	\label{def:matematica:logica-matematica:fundamentos:demonstracao}

	Uma demonstração é uma sequência finita de inferências lógicas que
	estabelece a validade de uma proposição a partir de axiomas,
	definições e proposições previamente demonstradas.
\end{defn}

\begin{axiom}[Princípio da Não Contradição]
	\label{axiom:matematica:logica-matematica:fundamentos:principio-nao-contradicao}

	Nenhuma proposição pode ser simultaneamente verdadeira e falsa.
\end{axiom}

\begin{axiom}[Princípio do Terceiro Excluído]
	\label{axiom:matematica:logica-matematica:fundamentos:principio-terceiro-excluido}

	Toda proposição é verdadeira ou falsa.
\end{axiom}

\begin{obs}
	Nesta obra não será feita distinção formal entre teoremas, lemas,
	corolários ou resultados equivalentes.

	Toda afirmação demonstrável será denominada proposição.
\end{obs}
\end{document}
